\section{Model formulation and implementation}
\label{sec:lavoro}

\subsection{General approach}
\label{sec:approccio}

The implementation combines two software packages with complementary roles.
OpenFOAM~\cite{openfoam13} provides the CFD framework, including mesh
management, spatial discretisation of the conservation equations and time
integration. The thermochemistry library \mpp~\cite{scoggins2020} provides
the modal energies, vibrational relaxation times and reaction rates, together
with the inversion from energies to temperatures. The solver advances the
energy equations, and the thermophysical model uses \mpp{} to recover the
two temperatures after each energy update.

The new components are compiled as \emph{dynamic libraries}, loaded at run
time and selected through the case configuration files. This architecture
keeps the project separate from the OpenFOAM installation and allows its
libraries to be compiled independently of the framework. It also confines
version-dependent changes to the interfaces with OpenFOAM, reducing the
maintenance required when updating the underlying framework.

\subsection{OpenFOAM}
\label{sec:openfoam}

\subsubsection{The framework and its one-temperature model}
\label{sec:openfoam-moduli}

OpenFOAM is an open-source CFD framework comprising
libraries and solvers based on the finite volume method. The governing
equations are discretised over a mesh of control volumes, with the primary
fields stored at cell centres. Its modular architecture supports additional
physical models without requiring changes to the core libraries.

The framework provides numerical methods relevant to the present application,
including central-upwind schemes for compressible flows with shock
waves~\cite{kurganov2001,greenshields2010}, species transport and chemical
kinetics. In the standard one-temperature formulation, a thermophysical model
relates the internal energy to temperature and supplies the equation of state
and transport properties. A separate chemistry model evaluates reaction rates
at the same temperature. This formulation therefore assumes the thermal
equilibrium described in Section~\ref{sec:una-temperatura}.

\subsubsection{Extensions to the framework}
\label{sec:openfoam-innesto}

The implementation extends the solver and the thermophysical model. The solver
derives from the module for high-speed compressible flows, which succeeds the
solver on which \hyfoam{} was based. The thermophysical solution stage is
modified to include the vibro-electronic energy equation and chemical source
terms, while retaining the existing numerical fluxes and the density,
momentum and pressure updates.

The thermophysical model extends the standard compressible-mixture model with
two additional fields, the vibro-electronic temperature and energy, and
provides the species, chemical-energy and vibro-electronic source terms. It
retains the standard OpenFOAM interface and is selected through the case
configuration. All calls to \mpp{} are confined to this model; the solver
accesses the library through the thermophysical interface.

\subsection{\mpp}
\label{sec:mutation}

\mpp{} supplies the gas properties and chemical reaction rates needed by the
solver. Its capabilities and integration with the solver are described below.

\subsubsection{Purpose and capabilities}
\label{sec:mutation-cosa}

\mpp~\cite{scoggins2020} is an open-source C++ library developed at the von
Karman Institute for Fluid Dynamics. It supplies thermodynamic and transport
properties, chemical reaction rates and energy-transfer source terms for use
in flow solvers.

A gas model is specified by its constituent species, thermodynamic database,
reaction mechanism and number of temperatures. The two-temperature
formulation used here distinguishes translational-rotational and
vibro-electronic energy. The thermodynamic state can be prescribed through
species densities and temperatures, species densities and energies, or mass
fractions, pressure and temperatures. The library evaluates the corresponding
properties and, when energies are prescribed, recovers the temperatures by
inversion.

\subsubsection{Coupling with the solver}
\label{sec:mutation-collegamento}

\mpp{} is compiled separately from OpenFOAM and linked to both the
thermophysical model and the zero-dimensional verification programs.

At initialisation, the thermophysical model reads the mixture definition,
number of temperatures, thermodynamic database and reaction mechanism from the
case configuration. It then constructs the corresponding \mpp{} mixture and
maps the OpenFOAM species to the library species by name. A missing species
causes initialisation to terminate with an error.

Following each energy update, the model supplies the partial densities
$\rho_s$, the internal energy per unit volume including the chemical energy
of formation, and the vibro-electronic energy per unit volume to \mpp{}.
The library uses these quantities to recover $\Ttr$ and $\Tve$. The
thermophysical model also evaluates the pressure-density relation, species
production rates $\dot\omega_s$ and energy-transfer source terms. The library
state is reset before each property or source-term evaluation to ensure
consistency with the current cell conditions.

\subsection{General thermochemical equations}
\label{sec:equazioni-generali}

For a spatially varying flow, a two-temperature model requires species and
energy balances containing both transport and local source terms. This
section introduces that general setting; the heat-bath equations implemented
and verified in this work are presented separately in
Section~\ref{sec:implementato}. The equations below describe the thermochemical
part of the problem and accompany the flow mass and momentum equations.

Here $t$ is time, $\rho$ is the gas density, $\mathbf{U}$ the velocity, $Y_s$ the mass
fraction of species $s$ and $\rho_s = \rho Y_s$ its partial density. The two
temperatures are $\Ttr$ and $\Tve$. The specific sensible internal energy $e$
excludes chemical formation energy and includes the vibro-electronic
component $\eve$.

\subsubsection{Species balance}

The conservation equation for each species is
\begin{equation}
  \frac{\partial (\rho Y_s)}{\partial t}
  + \nabla\cdot(\rho\,\mathbf{U}\,Y_s)
  + \nabla\cdot\mathbf{j}_s = \dot\omega_s ,
\end{equation}
where $\mathbf{j}_s$ is the diffusive mass flux and $\dot\omega_s$ the mass
production rate per unit volume. The flux terms describe transport between
neighbouring regions; the source describes local changes due to reactions.

\subsubsection{Sensible internal energy balance}

In the energy formulation of the underlying compressible-flow framework,
the sensible-energy balance with a chemical source is written as
\begin{equation}
  \frac{\partial (\rho e)}{\partial t}
  + \nabla\cdot\boldsymbol\phi_p
  + \frac{\partial (\rho K)}{\partial t}
  = Q_{ch}
  + \nabla\cdot(\tau\cdot\mathbf{U})
  + \nabla\cdot(\kappa\nabla\Ttr) ,
  \qquad
  Q_{ch} = -\sum_s h_{f,s}\,\dot\omega_s ,
\end{equation}
where $\boldsymbol\phi_p$ is the convective energy flux including pressure
work, $K = |\mathbf{U}|^2/2$ the specific kinetic energy, $\tau$ the viscous
stress, $\kappa$ the thermal conductivity and $h_{f,s}$ the specific enthalpy
of formation of species $s$. The source $Q_{ch}$ converts chemical formation
energy into sensible energy or vice versa. Since $e$ includes both thermal
modes, exchange between them redistributes energy internally and introduces
no additional source in this balance.

\subsubsection{Vibro-electronic energy balance}

A spatially varying two-temperature model also requires transport of
vibro-electronic energy. With the source terms considered here, its balance
can be expressed schematically as
\begin{equation}
  \frac{\partial (\rho\eve)}{\partial t}
  + \nabla\cdot(\rho\mathbf{U}\eve)
  + \nabla\cdot\mathbf{J}_{ve} = Q_{VT} + Q_{CV} ,
  \label{eq:eve-general}
\end{equation}
where $\mathbf{J}_{ve}$ groups the non-convective vibro-electronic energy
fluxes, including conduction and transport by species diffusion. The sources
$Q_{VT}$ and $Q_{CV}$ represent vibration--translation exchange and the
vibro-electronic energy associated with reactions, respectively. Although
Eq.~\eqref{eq:eve-general} represents the general spatially varying case,
the present verification considers spatially uniform heat baths, for which
all transport terms vanish and only the local sources $Q_{VT}$ and $Q_{CV}$
remain.

\subsection{Implemented heat-bath model and solution procedure}
\label{sec:implementato}

The cases verified in this work are closed, isolated heat baths represented
in OpenFOAM by a single cell. The gas is at rest and spatially uniform, so
$\mathbf{U}=0$ and all spatial gradients vanish. Consequently, convection,
diffusion, viscous work and heat conduction make no contribution. For the
two-temperature cases, the equations advanced in time reduce to
\begin{equation}
  \frac{\mathrm{d}(\rho Y_s)}{\mathrm{d}t} = \dot\omega_s ,
  \qquad
  \frac{\mathrm{d}(\rho e)}{\mathrm{d}t} = Q_{ch} ,
  \qquad
  \frac{\mathrm{d}(\rho\eve)}{\mathrm{d}t} = Q_{VT}+Q_{CV} .
\end{equation}
The total density remains constant in the closed, fixed volume. The
conserved internal energy includes the chemical formation energy; the
sensible part $e$ can change as the composition evolves.

OpenFOAM supplies the equation assembly and time integration, while the
thermophysical model uses \mpp{} to evaluate properties and source terms and
to recover temperatures. The inherited species and sensible-energy transport
operators remain in the solver but vanish in these cases. Transport in the
added vibro-electronic equation is not implemented and must be included
before applying the model to general flow configurations. The procedure
below therefore describes the thermochemical evolution verified under
heat-bath conditions.

\subsubsection{Initialisation: from temperatures to energies}
\label{sec:implementato-avvio}

The initial conditions specify the two temperatures, pressure and composition.
Before time integration begins, the thermophysical model obtains the
corresponding specific energies from \mpp{},
\begin{equation}
  e = \sum_s Y_s\big(e_s(\Ttr,\Tve) - h_{f,s}\big),
  \qquad
  \eve = \sum_s Y_s\,e_{ve,s}(\Tve) ,
\end{equation}
where $e_s$ is the specific internal energy of species $s$, including its
formation energy, $e_{ve,s}$ is its specific vibro-electronic energy and
$h_{f,s}$ its specific enthalpy of formation. Subtracting $h_{f,s}$ converts the \mpp{}
energy definition to the sensible-energy convention used by OpenFOAM. The
species and energy equations are then advanced in terms of $\rho Y_s$,
$\rho e$ and $\rho\eve$; the temperatures are recovered from the updated
state.

\subsubsection{Species and chemical-energy update}
\label{sec:implementato-specie}
\label{sec:implementato-energia}

The thermophysical update begins with the species equations. The model
obtains $\dot\omega_s$ from \mpp{} using the selected reaction mechanism;
in the single-cell heat bath these rates alone determine the evolution of
$\rho Y_s$. Nitrogen dissociation rates use Park's effective temperature $T_P$,
\begin{equation}
  T_P = \Ttr^{\,a}\,\Tve^{\,1-a}, \qquad a = 0.7 ,
\end{equation}
where $a$ is the exponent that weights the two temperatures. This choice
accounts for the reduction in dissociation rate when vibrational
excitation is below its equilibrium value at $\Ttr$. After the species
update, the mass fractions are normalised to sum to unity.

The sensible-energy equation uses the source
$Q_{ch}=-\sum_s h_{f,s}\dot\omega_s$ defined above. During dissociation,
sensible energy decreases as the chemical formation energy of the mixture
increases. The V--T exchange does not enter this equation because both
thermal modes are already included in $e$.

\subsubsection{Vibro-electronic source terms}
\label{sec:implementato-eve}

The implemented vibro-electronic update integrates the two local sources
$Q_{VT}$ and $Q_{CV}$.
The vibration--translation (V--T) source $Q_{VT}$ describes energy exchange
with the translational mode through relaxation towards
equilibrium~\cite{landau1936},
\begin{equation}
  Q_{VT} = \sum_m \rho_m\,
  \frac{e_{ve,m}(\Ttr) - e_{ve,m}(\Tve)}{\tau_m} ,
\end{equation}
where the sum extends over the molecular species and $\tau_m$ is the
relaxation time of species $m$, evaluated by \mpp{} using the Millikan--White
correlation~\cite{millikan1963} with Park's correction. The energy difference
provides the driving term, while $\tau_m$ determines the relaxation rate.
The source $Q_{CV}$ accounts for the vibro-electronic energy removed or added
by chemical reactions,
\begin{equation}
  Q_{CV} = \sum_s e_{ve,s}(\Tve)\,\dot\omega_s .
\end{equation}
\subsubsection{Recovering the temperatures}
\label{sec:implementato-temperature}

After the energy update, \mpp{} recovers the pair $(\Ttr,\Tve)$ that
reproduces the specified mixture energies,
\begin{equation}
  \Big(\rho_s,\; \rho\big(e + \textstyle\sum_s Y_s h_{f,s}\big),\; \rho\eve\Big)
  \;\longrightarrow\; (\Ttr,\,\Tve) .
\end{equation}
The inversion uses a Newton iteration based on the energy residuals and
corresponding heat capacities, terminating when the convergence tolerance is
satisfied.

The thermophysical model then evaluates the density at the current pressure
$p$, temperatures and composition, and updates the pressure-density
coefficient $\psi=\rho/p$ used by the pressure solution stage.

The same heat-bath equations are also integrated by the stand-alone
programs described in Section~\ref{sec:programmi0D}, with the case-specific
temperature formulations explained there. Comparison with those programs
checks the thermochemical coupling within OpenFOAM under these conditions;
it does not test spatial transport.

\subsection{Zero-dimensional verification programs}
\label{sec:programmi0D}

Four stand-alone programs reproduce the heat-bath case families of the
reference study without invoking OpenFOAM: pure nitrogen; molecular-atomic
nitrogen mixtures with and without reactions; nitrogen-oxygen mixtures; and
five-species air. Each program initialises the gas under the prescribed
conditions, evaluates the initial energies and advances the solution with a
time step of one nanosecond, as in the reference study. At each step, it
evaluates the source terms, advances the governing equations and recovers the
temperatures through \mpp{}. The total
internal energy, including chemical energy, remains constant in the closed,
isolated volume. The source-term functions are shared with the CFD solver to
ensure that both approaches use the same thermochemical formulation.

For non-reacting cases, the programs also calculate the expected equilibrium
temperature from conservation of the initial internal energy. This provides
an independent check of the final state reached by time integration.

Two case families require different temperature formulations. The
nitrogen-oxygen program evolves a separate vibrational energy and temperature
for each molecular species. It includes vibration--vibration (V--V) exchange
using the formulation of Knab \textit{et al.}~\cite{knab1995}, as adopted in
the reference study, with an exchange probability of $10^{-2}$. The
nitrogen-oxygen collision cross section, which is not reported in that study,
is taken from the \hyfoam{} data~\cite{hystrath}. The five-species air program
uses a single temperature, consistent with the corresponding reference case.

For the longest case, comprising one million time steps, the observed runtime
is approximately one second for the stand-alone program and almost half an
hour for the CFD solver. This difference reflects the additional operations
performed by the CFD framework. Even on a single-cell mesh, the solver
assembles and solves a matrix for each equation, performs two thermophysical
updates per time step and writes its output. The thermophysical model also
evaluates the library state separately for each species and source term,
whereas the stand-alone programs perform the updates through direct function
calls.

\subsection{Thermochemical data and reaction mechanisms}
\label{sec:dati}

The species data supplied to \mpp{} are taken from the reference study:
characteristic vibrational temperatures, enthalpies of formation and
electronic energy levels for the five air species. The vibrational relaxation
parameters are those distributed with \mpp{} and attributed to
Park~\cite{park1993}; for nitrogen and oxygen, they match the values used in
the reference study.

The treatment of electronic energy is specified explicitly in the reference
study only for nitrogen at 30\,000~K, where calculations with and without
electronic excitation are compared. Its treatment in the remaining cases is
not stated explicitly, although the discussion of molecular-atomic nitrogen
mentions the electronic contribution of atomic nitrogen. For the
molecular-atomic nitrogen and nitrogen-oxygen mixtures, the reported final
temperatures are reproduced only when excited electronic levels are excluded.
Including these levels gives a final temperature approximately one quarter
below the reference value for molecular-atomic nitrogen and more than one
half above it for nitrogen-oxygen. Excluding them reduces both discrepancies
to within 0.3\%. On this basis, two data sets are used: one includes the
electronic levels and is used for the nitrogen case with electronic
excitation; the other retains only the ground electronic level and is used
for the remaining cases. This choice is inferred from the reported results.

Nitrogen dissociation is represented by the irreversible reaction
$\mathrm{N_2 + N_2 \rightarrow 2\,N + N_2}$, using the Park rate constants
specified in the reference study. For five-species air, the reference
mechanism comprises nineteen irreversible reactions: fifteen dissociation
reactions and four exchange reactions. The corresponding Quantum-Kinetic
rate constants~\cite{scanlon2015} are cited but not tabulated in the study
and were obtained from the \hyfoam{} data~\cite{hystrath}. The air mechanism
distributed with \mpp{}, which uses Park rate constants, is also evaluated
for comparison.
